\section{USAMO 2008}
        \begin{enumerate}
        	\item (\textit{USAMO 2008}) (\textit{Titu Andreescu}) Prove that for each positive integer $n$, there are pairwise relatively prime integers $k_0,k_1\ldots,k_n$, all strictly greater than 1, such that $k_0k_1\cdots k_n-1$ is the product of two consecutive integers.


 



\textbf{Solution 1}

We will prove the problem for each nonnegative integer $n$.  We wish to show that
 \[k_0 k_1 \dotsc k_n = a^2+a+1,\]for some integer $a$.  We induct on $n$.  For our base case, $n=0$, we may let $a$ be positive integer.

 For the inductive step, suppose that $k_0, \dotsc, k_{n-1}$ are pairwise relatively prime integers such that
 \[k_0 \dotsm k_{n-1} = a^2+a+1,\]for some integer $a$.  Let $k_n = a^2 + 3a+3$.  Evidently, $k_n >1$.  Also,
 \[\gcd(a^2+a+1, k_n) = \gcd\bigl(a^2+a+1, k_n - (a^2+a+1) \bigr) = \gcd\bigl(a^2 + a+1, 2(a+1) \bigr) .\]Since $a^2+a+1 = a(a+1)+1$ is odd and relatively prime to $a+1$, it follows that $a^2+a+1$ and $k_n$ are relatively prime, so $k_n$ is relatively prime to each of $k_0, \dotsc, k_{n-1}$.  Finally,
 \begin{align*} k_0 k_1 \dotsm k_n &= (a^2+a+1)(a^2+3a+3)\\ &= \bigl[ (a+1)^2-a \bigr] \cdot \bigl[ (a+1)^2 + a+2 \bigr] \\ &= (a+1)^4 + 2(a+1)^2 - a^2-2a \\ &= (a+1)^4 + (a+1)^2 + 1 . \end{align*}
This completes the induction.  

 



\textbf{Solution 2}

\textbf{Lemma.}  If $p$ is prime such that $p\equiv 1 \pmod{3}$, there exists a residue $r$ such that $p \mid r^2+r+1$.

 \textit{Proof.} Let $g$ be a multiplicative generator of the nonzero integers mod p.  Set $r \equiv g^{(p-1)/3}$.  Then $r-1 \not\equiv 0 \pmod{p}$, but $r^3-1 \equiv 0 \pmod{p}$, so $\frac{r^3-1}{r-1} \equiv r^2 + r+1 \equiv 0 \pmod{p}$. 

 \textbf{End Lemma}

 By \href{https://artofproblemsolving.com/wiki/index.php?title=Dirichlet\%27s\_Theorem}{Dirichlet's Theorem}, there are infinitely many primes congruent to 1 (mod 3).  Let $p_0, \dotsc, p_n$ be $(n+1)$ such primes, and let $r_0, \dotsc, r_n$ be respective residues as described in the lemma.  By the \href{https://artofproblemsolving.com/wiki/index.php?title=Chinese\_Remainder\_Theorem}{Chinese Remainder Theorem}, there is a positive integer $a$ that satisfies the relation
 \[a \equiv r_i \pmod{p_i}\]for each integer $i \in [0,n]$.  Then
 \[p_0 \dotsc p_n \mid a^2+a+1 .\]Now, for $1 \le i \le n$, take $k_i$ to be the greatest power of $p_i$ that divides $a^2 + a +1$, and let $k_0 = (a^2+a+1)/(k_1 \dotsm k_n)$.  Since all the $k_i$ are pairwise relatively prime and are greater than 1, we are done. 

 



\textbf{Solution 3}

Firstly, we see that there are $n + 1$ numbers $k_i$. Since $n\geq1$, there are at least 2 values of $k_i$.Define a relatively prime partition to be a set of relatively prime numbers such that their product is equal to some natural number.

 Define $P$ to be the greatest possible cardinality of a relatively prime partition for that number.

 \textbf{Lemma 1.} All cardinalities of the relatively prime partitions of a number up to $P$ can be attained.

 \textit{Proof.} $P = P - 0$ satisfies the properties. For any $P - n$ which satisfies the properties, we can take any of the 2 $P - n$ numbers and multiply them together. Because they are both relatively prime to all the other numbers, their product is relatively prime to all the other numbers as well, and that results in $P - n - 1$ numbers which satisfy the conditions, unless $P - n = 1$, because there is only one number left. Therefore, all numbers of numbers of relatively prime factors from $1$ to $P - 1$ are attainable, if $P$ is attainable as well.

 \textbf{End Lemma}

 \textbf{Lemma 2.} $a^2 + a + 1$ can have arbitrarily many prime factors for some value of $a$.

 \textit{Proof.} Let $f(x) = x^2 + x + 1$. Let $a^2 + a + 1 = k$. Then $f(a + k) = a^2 + 2ak + k^2 + a + k + 1 = 2ak + k^2 + 2k$.After factoring, we get $k(k + 2 + 2a) = (a^2 + a + 1)(a^2 + 3a + 3)$.$a^2 + 3a + 3 \geq a^2 + a + 1$, because $a$ is strictly positive.Let their GCD=g.$g|2a + 2$$g|2a^2 + 2a$$g|a^2 - 3a - 3 + a^2 + 3a + 3, g|2a^2$$g|2a$, and so $g|2$, but $g|a^2 + a + 3$ which is strictly odd, so g=1.$(a^2 + a + 1)(a^2 + 3a + 3)$ therefore must contain prime factors not in $a^2 + a + 1$. Upon repeating this an arbitrary number of times, we get a number of the form $a^2 + a + 1$ which has arbitrarily many distinct prime factors.

 \textbf{End Lemma}

 Then $a(a + 1) + 1 = a^2 + a + 1$ would not have $n + 1 = P$ elements in the relatively prime partition for some value of n, and any value of a.It is possible to choose a value of $a$ such that $a^2 + a + 1$ has arbitrarily many unique prime factors, by our second Lemma, and so it is possible for $P$ to be arbitrarily high. By our first Lemma, all numbers up to P are possible values for the cardinality of some relatively prime partition, and so there always exists some number that has an arbitrary number of elements in a relatively prime partition. Because $a^2 + a + 1 = a(a + 1) + 1$, $a(a + 1)$ is the product of 2 consecutive integers, we see that the given statement is true because if $k_0k_1...k_n = a(a + 1) + 1$, then $k_0k_1...k_n - 1 = a(a + 1)$.

 



\textbf{Solution 4}

Write the relation to be proved as
 \[4k_0k_1\cdots k_n = 4a(a + 1) + 4 = (2a + 1)^2 + 3.\]There are infinitely many primes for which $-3$ is a quadratic residue. Let $2 < p_0 < p_1 < \cdots < p_n$ be such primes. Using the Chinese Remainder Theorem to specify $a$ modulo $p_n$, we can find an integer $a$ such that $(2a + 1)^2 + 3 = 4p_0p_1\cdots p_nm$ for some positive integer $m$. Grouping the factors of $m$ appropriately with the $p_i$'s, we obtain $(2a + 1)^2 + 3 = 4k_0k_1\cdots k_n$ with $k_i$ pairwise relatively prime. We then have $k_0k_1\cdots k_n - 1 = a(a + 1)$, as desired.

 



\textbf{Solution 5}

We are supposed to show that for every positive integer $n$, there is a positive integer $x$ such that $x(x + 1) + 1 = x^2 + x + 1$ has at least $n$ distinct prime divisors. We can actually prove a more general statement.

 \textbf{Claim.} Let $P(x) = a_dx^d + \cdots + a_1x + 1$ be a polynomial of degree $d\geq 1$ with integer coefficients. Then for every positive integer $n$, there is a positive integer $x$ such that $P(x)$ has at least $n$ distinct prime divisors.

 The proof follows from the following two lemmas.

 \textbf{Lemma 1.} The set
 \[Q = \{p\mid p\text{ a prime for which there is an integer }x\text{ such that }p\text{ divides }P(x)\}\]is infinite.

 \textit{Proof.} The proof is analogous to Euclid's proof that there are infinitely many primes. Namely, if we assume that there are only finitely many primes $p_1, p_2, \ldots, p_k$ in $Q$, then for each integer $m$, $P(mp_1p_2\cdots p_k)$ is an integer with no prime factors, which must equal $1$ or $-1$. However, since $P$ has degree $d\geq 1$, it takes each of the values $1$ and $-1$ at most $d$ times, a contradiction.

 \textbf{End Lemma}

 \textbf{Lemma 2.} Let $p_1, p_2, \ldots, p_n$, $n\geq 1$ be primes in $Q$. Then there is a positive integer $x$ such that $P(x)$ is divisible by $p_1p_2\cdots p_n$.

 \textit{Proof.} For $i = 1, 2, \ldots, n$, since $p_i\in Q$ we can find an integer $c_i$ such that $P(x)$ is divisible by $p_i$ whenever $x\equiv c_i\pmod{p_i}$. By the Chinese Remainder Theorem, the system of $n$ congruences $x\equiv c_i\pmod{p_i}$, $i = 1, 2, \ldots, n$ has positive integer solutions. For every positive integer $x$ that solves this system, $P(x)$ is divisible by $p_1p_2\cdots p_n$.

 \textbf{End Lemma}

 



\textbf{Solution 6}

Very similar to some above solutions.

 Let $k_{0}k_{1}\dots k_{n}-1=a(a+1)$. We wish to show this equation has a solution $(a,k_{0},\dots,k_{n})\in\mathbb{Z}^{n+2}$ for any $n\in\mathbb{N}$. If $a^{2}+a+1$ has at least $n+1$ distinct prime factors, we can set $a^{2}+a+1=p_{1}^{q_{1}}p_{2}^{q_{2}}\dots p_{m}^{q_{m}}$ for some $m\geq n$ where $p_{1},\dots p_{m}$ are distinct primes. Then, we let $k_{i}=p_{i+1}^{q_{i+1}}$ for all $0\leq i\leq n-1$ and set $k_{n}=p_{n+1}^{q_{n+1}}\dots p_{m}^{q_{m}}$. Thus, it remains to show $a^{2}+a+1$ can have at least $n+1$ distinct prime factors. We will do this by showing it can have arbitrarily many.

 Using an induction-type argument (it isn't exactly induction though), we first show that $a^{2}+a+1$ can have $1$ prime factor by setting $a=2$. Now, assume $a^{2}+a+1$ has $n$ distinct prime factors. We wish to show there exists some $b$ such that $b^{2}+b+1$ has more than $n$ distinct prime factors (we need not show it has $n+1$, since we only wish to demonstrate that the number of prime factors can be arbitrarily large rather than any integer). Set $b=a^{2}$. Then
 \begin{align*} b^{2}+b+1&=(b+1)^{2}-b\\ &=(a^{2}+1)^{2}-a^{2}\\ &=(a^{2}+a+1)(a^{2}-a+1). \end{align*}
Suppose $p$ is a positive integer dividing both $a^{2}+a+1$ and $a^{2}-a+1$. Then $p\mid(a^{2}+a+1)-(a^{2}-a+1)=2a$, and because $(a^{2}+a+1)$ is always odd, we have $2\nmid p$ so $p\mid a$. But this yields $p\mid\gcd(a,a^{2}+a+1)=1$, so $p=1$ and thus $a^{2}+a+1$ and $a^{2}-a+1$ are coprime. Therefore, as one can make $a^{2}-a+1>1$, we see that $a^{2}-a+1$ has at least one prime factor which does not divide $a^{2}+a+1$ (which has $n$ distinct prime factors). Therefore, $b^{2}+b+1$ has more than $n$ distinct prime factors.

 <i>Alternate solutions are always welcome. If you have a different, elegant solution to this problem, please \href{https://artofproblemsolving.com/wiki/index.php?title=2008\_USAMO\_Problems/Problem\_1&action=edit}{add it} to this page.</i>

 


	\item (\textit{USAMO 2008}) (\textit{Zuming Feng}) Let $ABC$ be an acute, \href{https://artofproblemsolving.com/wiki/index.php?title=Scalene}{scalene} triangle, and let $M$, $N$, and $P$ be the midpoints of $\overline{BC}$, $\overline{CA}$, and $\overline{AB}$, respectively. Let the \href{https://artofproblemsolving.com/wiki/index.php?title=Perpendicular}{perpendicular} \href{https://artofproblemsolving.com/wiki/index.php?title=Bisect}{bisectors} of $\overline{AB}$ and $\overline{AC}$ intersect ray $AM$ in points $D$ and $E$ respectively, and let lines $BD$ and $CE$ intersect in point $F$, inside of triangle $ABC$. Prove that points $A$, $N$, $F$, and $P$ all lie on one circle.


 



\textbf{Solution 1 (synthetic)}

Without Loss of Generality, assume $AB >AC$. It is sufficient to prove that $\angle OFA = 90^{\circ}$, where $O$ is the circumcenter of $\triangle{ABC},$ as this would immediately prove that $A,P,O,F,N$ are concyclic.By applying the Menelaus' Theorem in the Triangle $\triangle BFC$ for the transversal $E,M,D$, we have (in magnitude)
 \[\frac{FE}{EC} \cdot \frac{CM}{MB} \cdot \frac{BD}{DF} = 1 \iff \frac{FE}{EC} = \frac{DF}{BD}\] Here, we used that $BM=MC$, as $M$ is the midpoint of $BC$. Now, since $EC =EA$ and $BD=DA$, we have
 \[\frac{FE}{EA} = \frac{DF}{DA} \iff \frac{DA}{AE} = \frac{DF}{FE} \iff AF \text{ bisects exterior } \angle EFD\]Now, note that $OE$ bisects the exterior $\angle FED$ and $OD$ bisects exterior $\angle FDE$, making $O$ the $F$-excentre of $\triangle FED$. This implies that $OF$ bisects interior $\angle EFD$, making $OF \perp AF$, as was required.

 



\textbf{Solution 2 (complex)}

Let $A=1,B=b,C=c$ where $b,c$ all lie on the unit circle. Then $O$ is 0. As noted earlier, $(FOBC)$ is cyclic. We will find the ghost point $F',$ the second intersection of $OBC$ and $ANP$.

 We know that these two circles already intersect at $O$ so we can reflect over the line between their centers. The center of $ANPO$ is the midpoint of $AO$ namely $\frac12$. With the tangent formula and then taking the midpoint, we find that the center of $OBC$ is $\frac{bc}{b+c}.$ Then we want to find the reflection of 0 over the line through $\frac12$ and $\frac{bc}{b+c}.$ Then we get

 
 \begin{align*} f' &= \frac{\left(\frac{bc}{b+c}-\overline{\frac{bc}{b+c}}\right)\div2}{(1/2)-\overline{\frac{bc}{b+c}}}\\ &= \frac{\frac{bc-1}{2(b+c)}}{\frac{b+c-2}{2(b+c)}}\\ &= \frac{bc-1}{b+c-2}. \end{align*}


 Now it remains to show $\angle F'BA=\angle ABM;$ the other angle equality would follow by symmetry.

 Then we get:
 \begin{align*} \frac{f'-b}{b-a}\div\frac{b-a}{a-m} &=\frac{\frac{bc-1}{b+c-2}-b}{b-1}\div\frac{b-1}{1-\frac{b+c}2}\\ &=\frac{\frac{bc-1-b(b+c-2)}{b+c-2}(2-b+c)}{2(b-1)^2}\\ &=\frac{b^2-2b+1}{2(b-1)^2}\\ &=\frac12. \end{align*}
Thus $\measuredangle F'BA=\measuredangle BAM,$ so $F'=F$ and we're done.

 



\textbf{Solution 3 (synthetic)}


                    \begin{figure}[H]
                        \centering
                        \includegraphics[width=0.5\textwidth]{Solution/USAMO_I_2008_2_Solution_Figure_1.png}
                    \end{figure}
                    

 \href{https://artofproblemsolving.com/wiki/index.php?title=Without\_loss\_of\_generality}{Without loss of generality} $AB < AC$. The intersection of $NE$ and $PD$ is $O$, the circumcenter of $\triangle ABC$. 

 Let $\angle BAM = y$ and $\angle CAM = z$. Note $D$ lies on the perpendicular bisector of $AB$, so $AD = BD$. So $\angle FBC = \angle B - \angle ABD = B - y$. Similarly, $\angle FCB = C - z$, so $\angle BFC = 180 - (B + C) + (y + z) = 2A$. Notice that $\angle BOC$ intercepts the minor arc $BC$ in the \href{https://artofproblemsolving.com/wiki/index.php?title=Circumcircle}{circumcircle} of $\triangle ABC$, which is double $\angle A$. Hence $\angle BFC = \angle BOC$, so $BFOC$ is cyclic. 

 \textbf{Lemma.} $\triangle FEO$ is directly similar to $\triangle NEM$

 \textit{Proof.}
 \[\angle OFE = \angle OFC = \angle OBC = \frac {1}{2}\cdot (180 - 2A) = 90 - A\]since $F$, $E$, $C$ are collinear, $BFOC$ is cyclic, and $OB = OC$. Also
 \[\angle ENM = 90 - \angle MNC = 90 - A\]because $NE\perp AC$, and $MNP$ is the medial triangle of $\triangle ABC$ so $AB \parallel MN$. Hence $\angle OFE = \angle ENM$.

 Notice that $\angle AEN = 90 - z = \angle CEN$ since $NE\perp AC$. $\angle FED = \angle MEC = 2z$. Then
 \[\angle FEO = \angle FED + \angle AEN = \angle CEM + \angle CEN = \angle NEM\]Hence $\angle FEO = \angle NEM$. 

 Hence $\triangle FEO$ is similar to $\triangle NEM$ by AA similarity. It is easy to see that they are oriented such that they are directly similar.

 \textbf{End Lemma}

 
                    \begin{figure}[H]
                        \centering
                        \includegraphics[width=0.5\textwidth]{Solution/USAMO_I_2008_2_Solution_Figure_2.png}
                    \end{figure}
                    

 By the similarity in the Lemma, $FE: EO = NE: EM\implies FE: EN = OE: EM$. $\angle FEN = \angle OEM$ so $\triangle FEN\sim\triangle OEM$ by SAS similarity. Hence
 \[\angle EMO = \angle ENF = \angle ONF\]Using essentially the same angle chasing, we can show that $\triangle PDM$ is directly similar to $\triangle FDO$. It follows that $\triangle PDF$ is directly similar to $\triangle MDO$. So
 \[\angle EMO = \angle DMO = \angle DPF = \angle OPF\]Hence $\angle OPF = \angle ONF$, so $FONP$ is cyclic. In other words, $F$ lies on the circumcircle of $\triangle PON$. Note that $\angle ONA = \angle OPA = 90$, so $APON$ is cyclic. In other words, $A$ lies on the circumcircle of $\triangle PON$. $A$, $P$, $N$, $O$, and $F$ all lie on the circumcircle of $\triangle PON$. Hence $A$, $P$, $F$, and $N$ lie on a circle, as desired.

 



\textbf{Solution 4 (synthetic)}

This solution utilizes the \textit{phantom point method.} Clearly, APON are cyclic because $\angle OPA = \angle ONA = 90$. Let the circumcircles of triangles $APN$ and $BOC$ intersect at $F'$ and $O$. 

 \textbf{Lemma.} If $A,B,C$ are points on circle $\omega$ with center $O$, and the tangents to $\omega$ at $B,C$ intersect at $Q$, then $AP$ is the symmedian from $A$ to $BC$.

 \textit{Proof.} This is fairly easy to prove (as H, O are isogonal conjugates, plus using SAS similarity), but the author lacks time to write it up fully, and will do so soon.

 \textbf{End Lemma}

 It is easy to see $Q$ (the intersection of ray $OM$ and the circumcircle of $\triangle BOC$) is colinear with $A$ and $F'$, and because line $OM$ is the diameter of that circle, $\angle QBO = \angle QCO = 90$, so $Q$ is the point $Q$ in the lemma; hence, we may apply the lemma. From here, it is simple angle-chasing to show that $F'$ satisfies the original construction for $F$, showing $F=F'$; we are done.

 


 



\textbf{Solution 5 (trigonometric)}

By the \href{https://artofproblemsolving.com/wiki/index.php?title=Law\_of\_Sines}{Law of Sines}, $\frac {\sin\angle BAM}{\sin\angle CAM} = \frac {\sin B}{\sin C} = \frac bc = \frac {b/AF}{c/AF} = \frac {\sin\angle AFC\cdot\sin\angle ABF}{\sin\angle ACF\cdot\sin\angle AFB}$. Since $\angle ABF = \angle ABD = \angle BAD = \angle BAM$ and similarly $\angle ACF = \angle CAM$, we cancel to get $\sin\angle AFC = \sin\angle AFB$. Obviously, $\angle AFB + \angle AFC > 180^\circ$ so $\angle AFC = \angle AFB$.

 Then $\angle FAB + \angle ABF = 180^\circ - \angle AFB = 180^\circ - \angle AFC = \angle FAC + \angle ACF$ and $\angle ABF + \angle ACF = \angle A = \angle FAB + \angle FAC$. Subtracting these two equations, $\angle FAB - \angle FCA = \angle FCA - \angle FAB$ so $\angle BAF = \angle ACF$. Therefore, $\triangle ABF\sim\triangle CAF$ (by AA similarity), so a spiral similarity centered at $F$ takes $B$ to $A$ and $A$ to $C$. Therefore, it takes the midpoint of $\overline{BA}$ to the midpoint of $\overline{AC}$, or $P$ to $N$. So $\angle APF = \angle CNF = 180^\circ - \angle ANF$ and $APFN$ is cyclic.

 



\textbf{Solution 6 (isogonal conjugates)}


                    \begin{figure}[H]
                        \centering
                        \includegraphics[width=0.5\textwidth]{Solution/USAMO_I_2008_2_Solution_Figure_3.png}
                    \end{figure}
                    

 Construct $T$ on $AM$ such that $\angle BCT = \angle ACF$. Then $\angle BCT = \angle CAM$. Then $\triangle AMC\sim\triangle CMT$, so $\frac {AM}{CM} = \frac {CM}{TM}$, or $\frac {AM}{BM} = \frac {BM}{TM}$. Then $\triangle AMB\sim\triangle BMT$, so $\angle CBT = \angle BAM = \angle FBA$. Then we have

 $\angle CBT = \angle ABF$ and $\angle BCT = \angle ACF$. So $T$ and $F$ are \href{https://artofproblemsolving.com/wiki/index.php?title=Isogonal\_conjugate}{isogonally conjugate}. Thus $\angle BAF = \angle CAM$. Then

 $\angle AFB = 180 - \angle ABF - \angle BAF = 180 - \angle BAM - \angle CAM = 180 - \angle BAC$.

 If $O$ is the \href{https://artofproblemsolving.com/wiki/index.php?title=Circumcenter}{circumcenter} of $\triangle ABC$ then $\angle BFC = 2\angle BAC = \angle BOC$ so $BFOC$ is cyclic. Then $\angle BFO = 180 - \angle BOC = 180 - (90 - \angle BAC) = 90 + \angle BAC$.

 Then $\angle AFO = 360 - \angle AFB - \angle BFO = 360 - (180 - \angle BAC) - (90 + \angle BAC) = 90$. Then $\triangle AFO$ is a right triangle.

 Now by the \href{https://artofproblemsolving.com/wiki/index.php?title=Homothety}{homothety} centered at $A$ with ratio $\frac {1}{2}$, $B$ is taken to $P$ and $C$ is taken to $N$. Thus $O$ is taken to the circumcenter of $\triangle APN$ and is the midpoint of $AO$, which is also the circumcenter of $\triangle AFO$, so $A,P,N,F,O$ all lie on a circle.

 



\textbf{Solution 7 (symmedians)}

Median $AM$ of a triangle $ABC$ implies $\frac {\sin{BAM}}{\sin{CAM}} = \frac {\sin{B}}{\sin{C}}$.Trig ceva for $F$ shows that $AF$ is a symmedian.Then $FP$ is a median, use the lemma again to show that $AFP = C$, and similarly $AFN = B$, so you're done.

 


 



\textbf{Solution 8 (inversion) (Official Solution #2)}

Invert the figure about a circle centered at $A$, and let $X'$ denote the image of the point $X$ under this inversion. Find point $F_1'$ so that $AB'F_1'C'$ is a parallelogram and let $Z'$ denote the center of this parallelogram. Note that $\triangle BAC\sim\triangle C'AB'$ and $\triangle BAD\sim\triangle D'AB'$. Because $M$ is the midpoint of $BC$ and $Z'$ is the midpoint of $B'C'$, we also have $\triangle BAM\sim\triangle C'AZ'$. Thus
 \[\angle AF_1'B' = \angle F_1'AC' = \angle Z'AC' = \angle MAB = \angle DAB = \angle DBA = \angle AD'B'.\]Hence quadrilateral $AB'D'F_1'$ is cyclic and, by a similar argument, quadrilateral $AC'E'F_1'$ is also cyclic. Because the images under the inversion of lines $BDF$ and $CFE$ are circles that intersect in $A$ and $F'$, it follows that $F_1' = F'$.

 Next note that $B'$, $Z'$, and $C'$ are collinear and are the images of $P'$, $F'$, and $N'$, respectively, under a homothety centered at $A$ and with ratio $1/2$. It follows that $P'$, $F'$, and $N'$ are collinear, and then that the points $A$, $P$, $F$, and $N$ lie on a circle.

 
                    \begin{figure}[H]
                        \centering
                        \includegraphics[width=0.5\textwidth]{Solution/USAMO_I_2008_2_Solution_Figure_4.png}
                    \end{figure}
                    

 



\textbf{Solution 9 (Official Solution #1)}

Let $O$ be the circumcenter of triangle $ABC$. We prove that
 \[\angle APO = \angle ANO = \angle AFO = 90^\circ.\qquad\qquad (1)\]It will then follow that $A, P, O, F, N$ lie on the circle with diameter $\overline{AO}$. Indeed, the fact that the first two angles in $(1)$ are right is immediate because $\overline{OP}$ and $\overline{ON}$ are the perpendicular bisectors of $\overline{AB}$ and $\overline{AC}$, respectively. Thus we need only prove that $\angle AFO = 90^\circ$.

 
                    \begin{figure}[H]
                        \centering
                        \includegraphics[width=0.5\textwidth]{Solution/USAMO_I_2008_2_Solution_Figure_5.png}
                    \end{figure}
                    

 We may assume, without loss of generality, that $AB > AC$. This leads to configurations similar to the ones shown above. The proof can be adapted to other configurations. Because $\overline{PO}$ is the perpendicular bisector of $\overline{AB}$, it follows that triangle $ADB$ is an isosceles triangle with $AD = BD$. Likewise, triangle $AEC$ is isosceles with $AE = CE$. Let $x = \angle ABD = \angle BAD$ and $y = \angle CAE = \angle ACE$, so $x + y = \angle BAC$.

 Applying the Law of Sines to triangles $ABM$ and $ACM$ gives
 \[\frac{BM}{\sin x} = \frac{AB}{\sin\angle BMA}\quad\text{and}\quad\frac{CM}{\sin y} = \frac{AC}{\sin\angle CMA}.\]Taking the quotient of the two equations and noting that $\sin\angle BMA = \sin\angle CMA$, we find
 \[\frac{BM}{CM}\frac{\sin y}{\sin x} = \frac{AB}{AC}\frac{\sin\angle CMA}{\sin\angle BMA} = \frac{AB}{AC}.\]Because $BM = MC$, we have
 \[\frac{\sin x}{\sin y} = \frac{AC}{AB}.\qquad\qquad (2)\]Applying the Law of Sines to triangles $ABF$ and $ACF$, we find
 \[\frac{AF}{\sin x} = \frac{AB}{\sin\angle AFB}\quad\text{and}\quad\frac{AF}{\sin y} = \frac{AC}{\sin\angle AFC}.\]Taking the quotient of the two equations yields
 \[\frac{\sin x}{\sin y} = \frac{AC}{AB}\frac{\sin\angle AFB}{\sin\angle AFC},\]so by $(2)$,
 \[\sin\angle AFB = \sin\angle AFC.\qquad\qquad (3)\]Because $\angle ADF$ is an exterior angle to triangle $ADB$, we have $\angle EDF = 2x$. Similarly, $\angle DEF = 2y$. Hence
 \[\angle EFD = 180^\circ - 2x - 2y = 180^\circ - 2\angle BAC.\]Thus $\angle BFC = 2\angle BAC = \angle BOC$, so $BOFC$ is cyclic. In addition,
 \[\angle AFB + \angle AFC = 360^\circ - 2\angle BAC > 180^\circ,\]and hence, from $(3)$, $\angle AFB = \angle AFC = 180^\circ - \angle BAC$. Because $BOFC$ is cyclic and $\triangle BOC$ is isosceles with vertex angle $\angle BOC = 2\angle BAC$, we have $\angle OFB = \angle OCB = 90^\circ - \angle BAC$. Therefore,
 \[\angle AFO = \angle AFB - \angle OFB = (180^\circ - \angle BAC) - (90^\circ - \angle BAC) = 90^\circ.\]This completes the proof.

 



\textbf{Solution 10 (Anti-Steiner point)}

First, let $O$ be the circumcenter of $\triangle{ABC}$. Note that since $AP \perp PO$ and $AN \perp NO$, then $A$, $N$, $O$ and $P$, are concyclic.

 Notice that $NM \parallel AB$ and $AB \perp PO$, which means $NM \perp PO$. It follows analogously that $PM \perp NO$. This means that $M$ is the orthocenter of triangle ${NOP}$.

 Now consider what happens when we reflect line $AM$ over line $PO$. Since $PO$ is just the perpendicular bisector of $AB$, point $A$ will map to point $B$, and point $D$, which is both line $AM$ and $PO$, will map to itself. Therefore, we get line $BD$ from this reflection. It follows analogously that by reflecting line $AM$ over line $NO$, we get $CE$.

 Since line $AM$ passes through $M$, the orthocenter of triangle ${NOP}$, its reflections over sides $PO$ and $NO$, which correspond to $BD$ and $CE$ respectively, intersect on $\odot{NOP}$ (this is the Anti-Steiner point application). Therefore, $F$, the intersection of $BD$ and $CE$, lies on $\odot{NOP}$.

 Since $A$, $N$, $O$, and $P$ are concyclic, then $\odot{NOP}$ is the same as $\odot{ANP}$, so $F$ lies on $\odot{ANP}$, which is the same as saying $A$, $N$, $F$, and $P$ are concyclic, as desired.

 



\textbf{Solution 11 (Bary from EGMO)}

First, we find the coordinates of $D$. As $D$ lies on $AM$, we know $D=(t:1:1)$ for some $T$. Now, we find 
 \[0=b^2(t-1)+(a^2-c^2) \implies t=\frac{c^2+b^2-a^2}{b^2}.\]

 Thus we obtain 
 \[D = (2S_A:c^2:c^2).\]

 (Note we are using Conway's notation here). Similarly, $E = (2S_A:b^2:b^2)$, and it follows that 
 \[F=(2S_A:b^2:c^2).\]

 The sum of the coordinates of $F$ is 
 \[(b^2+c^2-a^2)+b^2+c^2=2b^2+2c^2-a^2.\]

 Hence, the reflection of $A$ over $F$ is simply 
 \[2F-A=(-a^2:2b^2:2c^2).\]

 Then, it is obvious that $F$ lies on $(ABC):-a^2yz-b^2zx-c^2xy=0$.

 <i>Alternate solutions are always welcome. If you have a different, elegant solution to this problem, please \href{https://artofproblemsolving.com/wiki/index.php?title=2008\_USAMO\_Problems/Problem\_2&action=edit}{add it} to this page.</i>

 


	\item (\textit{USAMO 2008}) (\textit{Gabriel Carroll}) Let $n$ be a positive integer. Denote by $S_n$ the set of points $(x, y)$ with integer coordinates such that

 \[\left|x\right|+\left|y+\frac{1}{2}\right|<n\]
A path is a sequence of distinct points $(x_1,y_1),(x_2,y_2),\ldots,(x_\ell,y_\ell)$ in $S_n$ such that, for $i=2,\ldots,\ell$, the distance between $(x_i,y_i)$ and $(x_{i-1},y_{i-1})$ is $1$ (in other words, the points $(x_i,y_i)$ and $(x_{i-1},y_{i-1})$ are neighbors in the lattice of points with integer coordinates). Prove that the points in $S_n$ cannot be partitioned into fewer than $n$ paths (a partition of $S_n$ into $m$ paths is a set $\mathcal{P}$ of $m$ nonempty paths such that each point in $S_n$ appears in exactly one of the $m$ paths in $\mathcal{P}$).


 



\textbf{Solution 1}

Color all the points in $S_{n}$ red or black such that each row of points starts and ends with a red point, and alternates between red and black for each point in the row. This creates a checkerboard pattern, except that the middle two rows are identical.  

 
                    \begin{figure}[H]
                        \centering
                        \includegraphics[width=0.5\textwidth]{Solution/USAMO_I_2008_3_Solution_Figure_1.png}
                    \end{figure}
                             
                    \begin{figure}[H]
                        \centering
                        \includegraphics[width=0.5\textwidth]{Solution/USAMO_I_2008_3_Solution_Figure_2.png}
                    \end{figure}
                    Examples for $n=4$

 Suppose there is a partition of $m < n$ paths that works. For each of the $m$ paths, we break a path into two separate paths wherever there are consecutive red points. This only happens only with the $n$ red pairs in the middle two rows. We now have $m + l,\, l \le n$ paths whose colors are strictly alternating. In a path with alternating colors, there can be at most one more red dot than the number of black dots (which occurs when both endpoints are red). Thus there can be no more than $m + l$ red dots than black dots in $S_n$. 

 There are $2n$ more red points than black points, so
 \[2n \le m + l \le m + n\]which implies that $n \le m$, contradiction.   

 



\textbf{Solution 2}

Suppose you have a partition of less than $n$ paths. Then start from $(0,n - 1)$ and work your way down to $(0, - n)$ along the right edge of $S_n$ performing the following algorithm when going from each point.

 
                    \begin{figure}[H]
                        \centering
                        \includegraphics[width=0.5\textwidth]{Solution/USAMO_I_2008_3_Solution_Figure_3.png}
                    \end{figure}
                    Example for $S_4$

 If a path already has the point you're coming from and the point you're going to as adjacent in the path, do nothing.

 If not the first case, and the point you're coming from and the point you're going to are both endpoints of paths, join them (one less path).

 If not the first case, and there is exactly one endpoint among the point you're coming from and the point you're going to, remove a segment from the one in the middle of a path (paths +1) and now you have two endpoints; join them (paths -1). If the point in the middle of the path is the one you're coming from, remove the segment that doesn't come from the point that came before it. The number of paths remains the same.

 If not the first case, and both points are in the middle of a path, which is impossible. Since the two points are in the middle of the path, they are adjacent to two other points. But they are also adjacent to each other, so both are adjacent to at least three points. Then, no two points adjacent to at least three others are travelled consecutively.

 Thus, given any configuration we can create the path shown above without increasing the number of paths. The remaining points form $S_{n - 1}$. Now suppose for some positive integer $n$, we have a partition into less than $n$ paths. Apply the algorithm to make the path shown above. Then the remaining points are partitioned into less than $n - 1$ paths. Apply the algorithm again. Repeat until we show that $S_1$ is partitioned into less than $1$ path, contradiction.

 



\textbf{Solution 3}

We induct on $n$ to prove that an $n - 1$-path partition is impossible. This can be restated as saying that we cannot fit $2n^2 - n + 1$ edges into the graph with each vertex having degree at most 2. The base case is trivial (2 edges when there only is 1). Suppose it was possible for $n + 1$, and we'll prove it's possible for $n$.

 Call the 'complete' network (graph) of points $G_{n + 1}$ and the $n$-path partition for $n + 1$ is a subgraph of $G_{n + 1}$, $H_{n + 1}$. There are a total of $2(n + 1)^2 - n = 2n^2 + 3n + 2$ edges in $H_{n + 1}$ out of $4n^2 + 4n + 1$ in $G_{n + 1}$ (not hard to prove), and each vertex has maximum order 2. Now consider the complement of $H_{n + 1}$ in $G_{n + 1}$, that is, take all the edges in $G_{n + 1}$ that are not in $H_{n + 1}$ and of course all the vertices. Call this $F_{n + 1}$ We will trim $F_{n + 1}$ down a bit to get $H_n$.

 First, call the points in $S_{n + 1}$ but not $S_n$ exterior points, and the points in $S_n$ interior points. $F_{n + 1}$ has $2n^2 + n - 1$ edges. Since the interior points have degree 4 in $G_{n + 1}$, they have degree at least 2 in $F_{n + 1}$, for a total degree of $4n^2 + m$ where $0\le m\le 2n - 2$ since the total overall degree is $2(2n^2 + n - 1) = 4n^2 + 2n - 2$. This makes the total degree of the exterior points $2n - 2 - m$. Since the exterior points are not connected in $G_{n + 1}$, these $2n - 2 - m$ edges must connect on exterior point to an interior point.

 I then take away $2n - 2$ edges and the exterior points, yielding $2n^2 - n + 1$ edges and the interior points, which is $H_n$, as follows: First take away those $2n - 2 - m$ interior-exterior edges and shave off all the exterior points. Then the total degree (now just of the interior points) must be $4n^2 - 2n + 2 + 2m$.

 Let the number of 0-degree vertices be $a$, 1-degree vertices be $b$, 2-degree vertices be $c$, 3-degree vertices be $d$, and 4-degree vertices be $e$. Since all the vertices began as 2-degree or higher, $2a + b\le2n - 2 - m$. Moreover, the total degree is $b + 2c + 3d + 4e = 2(2n^2) + (d + 2e) - (b + 2a) = 4n^2 - 2n + 2 + 2m$. So $d + 2e = b + 2a - (2n - 2 + 2m)\le (2n - 2 - m) - (2n - 2 - 2m) = m$ So we can use our last $m$ edges to remove any excess degree on vertices with degree 3 or 4, yielding $H_n$.

 


 <i>Alternate solutions are always welcome. If you have a different, elegant solution to this problem, please \href{https://artofproblemsolving.com/wiki/index.php?title=2008\_USAMO\_Problems/Problem\_3&action=edit}{add it} to this page.</i>

 


	\item (\textit{USAMO 2008}) (\textit{Gregory Galparin}) Let $\mathcal{P}$ be a \href{https://artofproblemsolving.com/wiki/index.php?title=Convex\_polygon}{convex polygon} with $n$ sides, $n\ge3$. Any set of $n-3$ diagonals of $\mathcal{P}$ that do not intersect in the interior of the polygon determine a \textit{triangulation} of $\mathcal{P}$ into $n - 2$ triangles. If $\mathcal{P}$ is regular and there is a triangulation of $\mathcal{P}$ consisting of only isosceles triangles, find all the possible values of $n$.


 



\textbf{Solution}

We label the vertices of $\mathcal{P}$ as $P_0, P_1, P_2, \ldots, P_n$. Consider a diagonal $d = \overline{P_a\,P_{a+k}},\,k \le n/2$ in the triangulation. We show that $k$ must have the form $2^{m}$ for some nonnegative integer $m$.

 This diagonal \href{https://artofproblemsolving.com/wiki/index.php?title=Partition}{partitions} $\mathcal{P}$ into two regions $\mathcal{Q},\, \mathcal{R}$, and is the side of an isosceles triangle in both regions. \href{https://artofproblemsolving.com/wiki/index.php?title=Without\_loss\_of\_generality}{Without loss of generality} suppose the area of $Q$ is less than the area of $R$ (so the \href{https://artofproblemsolving.com/wiki/index.php?title=Circumcenter}{center} of $P$ does not lie in the interior of $Q$); it follows that the lengths of the edges and diagonals in $Q$ are all smaller than $d$. Thus $d$ must the be the base of the isosceles triangle in $Q$, from which it follows that the isosceles triangle is $\triangle P_aP_{a+k/2}\,P_{a+k}$, and so $2|k$. Repeating this process on the legs of isosceles triangle ($\overline{P_aP_{a+k/2}},\,\overline{P_{a+k}P_{a+k/2}}$), it follows that $k = 2^m$ for some positive integer $m$ (if we allow \href{https://artofproblemsolving.com/wiki/index.php?title=Degenerate}{degeneracy}, then we can also let $m=0$). 

 <table><tbody><tr><td>
                    \begin{figure}[H]
                        \centering
                        \includegraphics[width=0.5\textwidth]{Solution/USAMO_I_2008_4_Solution_Figure_1.png}
                    \end{figure}
                          </td><td>
                    \begin{figure}[H]
                        \centering
                        \includegraphics[width=0.5\textwidth]{Solution/USAMO_I_2008_4_Solution_Figure_2.png}
                    \end{figure}
                    </td></tr><tr><td>An example for $n=17$,
$k=8$</td><td>An isosceles triangle containing
 the center for 
 $n=20$, $(x,y,z) = (0,8,16)$</td></tr></tbody></table>

 Now take the isosceles triangle $P_xP_yP_z,\,0 \le x < y < z < n$ in the triangulation that contains the center of $\mathcal{P}$ in its interior; if a diagonal passes through the center, select either of the isosceles triangles with that diagonal as an edge. Without loss of generality, suppose $P_xP_y = P_yP_z$. From our previous result, it follows that there are $2^a$ edges of $P$ on the minor arcs of $P_xP_y,\, P_yP_z$ and $2^b$ edges of $P$ on the minor arc of $P_zP_x$, for positive integers $a,\,b$. Therefore, we can write
 \[n = 2 \cdot 2^a + 2^b = 2^{a+1} + 2^{b},\] so $n$ must be the sum of two powers of $2$. 

 
We now claim that this condition is sufficient. Suppose without loss of generality that $a+1 \ge b$; then we rewrite this as 
 \[n = 2^{b}(2^{a-b+1}+1).\]

 
                \begin{itemize}
                    \item \textit{Lemma 1}: All regular polygons with $n = 2^k + 1$ or $n=4$ have triangulations that meet the conditions.
                \end{itemize}


                By \href{https://artofproblemsolving.com/wiki/index.php?title=Induction}{induction}, it follows that we can cover all the desired $n$.

 For $n = 3,4$, this is trivial. For $k>1$, we construct the diagonals of equal length $\overline{P_0P_{2^{k-1}}}$ and $\overline{P_{2^{k-1}+1}P_0}$. This partitions $\mathcal{P}$ into $3$ regions: an isosceles $\triangle P_0P_{2^{k-1}}P_{2^{k-1}+1}$, and two other regions. For these two regions, we can recursively construct the isosceles triangles defined above in the second paragraph. It follows that we have constructed $2(2^{k-1}-1) + (1) = 2^k-1 = n-2$ isosceles triangles with non-intersecting diagonals, as desired.  

 
                    \begin{figure}[H]
                        \centering
                        \includegraphics[width=0.5\textwidth]{Solution/USAMO_I_2008_4_Solution_Figure_3.png}
                    \end{figure}
                    
An example for $n=17 = 2^{4}+1$

 
                \begin{itemize}
                    \item \textit{Lemma 2}: If a regular polygon with $n$ sides has a working triangulation, then the regular polygon with $2n$ sides also has a triangulation that meets the conditions.
                \end{itemize}


                We construct the diagonals $\overline{P_0P_2},\ \overline{P_2P_4},\ \ldots \overline{P_{2n-2}P_0}$. This partitions $\mathcal{P}$ into $n$ isosceles triangles of the form $\triangle P_{2k}P_{2k+1}P_{2k+2}$, as well as a central regular polygon with $n$ sides. However, we know that there exists a triangulation for the $n$-sided polygon that yields $n-2$ isosceles triangles. Thus, we have created $(n) + (n-2) = 2n-2$ isosceles triangles with non-intersecting diagonals, as desired. 

 
                    \begin{figure}[H]
                        \centering
                        \includegraphics[width=0.5\textwidth]{Solution/USAMO_I_2008_4_Solution_Figure_4.png}
                    \end{figure}
                    
An example for $n=10,\, n/2 = 5$

 In summary, the answer is all $n$ that can be written in the form $2^{a+1} + 2^{b},\, a,b \ge 0$. Alternatively, this condition can be expressed as either $n=2^{k},\, k \ge 2$ (this is the case when $a+1 = b$) or $n$ is the sum of two distinct powers of $2$, where $1= 2^0$ is considered a power of $2$.

 


	\item (\textit{USAMO 2008}) (\textit{Kiran Kedlaya}) Three nonnegative real numbers $r_1$, $r_2$, $r_3$ are written on a blackboard. These numbers have the property that there exist integers $a_1$, $a_2$, $a_3$, not all zero, satisfying $a_1r_1+a_2r_2+a_3r_3=0$. We are permitted to perform the following operation: find two numbers $x$, $y$ on the blackboard with $x\le y$, then erase $y$ and write $y-x$ in its place. Prove that after a finite number of such operations, we can end up with at least one $0$ on the blackboard.


 



\textbf{Solution}

Every time we perform an operation on the numbers on the blackboard $R = \left < r_1, r_2, r_3 \right >$, we perform the corresponding operation on the integers $A = \left < a_1, a_2, a_3 \right >$ so that $R \cdot A = 0$ continues to hold.  (For example, if we replace $r_1$ with $r_1 - r_2$ then we replace $a_2$ with $a_1 + a_2$.)

 It's possible to show we can always pick an operation so that $|A|^2$ is strictly decreasing. \href{https://artofproblemsolving.com/wiki/index.php?title=Without\_loss\_of\_generality}{Without loss of generality}, let $r_3 > r_2 > r_1$ and $a_3$ be positive. Then it cannot be true that both $a_1$ and $a_2$ are at least $\frac { - a_3}{2}$, or else $a_1r_1 + a_2r_2 + a_3r_3 > 0$. Without loss of generality, let $a_1 < \frac { - a_3}{2}$. Then we can replace $a_1$ with $a_1 + a_3$ and $r_3$ with $r_3 - r_1$ to make $|A|$ smaller. Since it is a strictly decreasing sequence of positive integers, after a finite number of operations we have $a_3 = 0$. We can now  see that this result holds for $(r_1,r_2,0)$ if and only if it holds for $(1,\frac{r_2}{r_1},0)$. We can see that $\frac{r_2}{r_1}$ is a rational number given that $a_3$ = 0. It is a well known result of the euclidean algorithm that if we continue to perform these operations, $r_1$ or $r_2$ will eventually be 0.

 


 


 


	\item (\textit{USAMO 2008}) (<i>\href{https://artofproblemsolving.com/wiki/index.php?title=Sam\_Vandervelde}{Sam Vandervelde}</i>) At a certain mathematical conference, every pair of mathematicians are either friends or strangers. At mealtime, every participant eats in one of two large dining rooms. Each mathematician insists upon eating in a room which contains an even number of his or her friends. Prove that the number of ways that the mathematicians may be split between the two rooms is a power of two (i.e., is of the form $2^k$ for some positive integer $k$).


 



\textbf{Solution 1 (linear algebra)}

Make the obvious re-interpretation as a \href{https://artofproblemsolving.com/wiki/index.php?title=Graph}{graph}. Let $f : G \to \{0,1\}$ be an \href{https://artofproblemsolving.com/wiki/index.php?title=Indicator\_function&action=edit&redlink=1}{indicator function} with $f(v) = 0$ if a vertex is in the first \href{https://artofproblemsolving.com/wiki/index.php?title=Partition}{partition} and $f(v) = 1$ otherwise (this corresponds, in the actual problem, to putting a mathematician in the first or second room). Then look at $f$ as a function into the \href{https://artofproblemsolving.com/wiki/index.php?title=Field}{field} with two elements, $F_2$. Let $V$ be the \href{https://artofproblemsolving.com/wiki/index.php?title=Vector\_space}{vector space} of all such functions. Define the linear operator $A : V \to V$ as

 
 \[(Af)(v) = \sum_{v \sim w} f(v) - f(w)\]

 where $\sim$ denotes adjacency. (Note that we can also think of $A$ as a matrix, which is essentially the adjacency matrix where the diagonal is changed to be $1$ whenever the degree is odd; in more technical terms, it is the Laplacian of $G$ over $F_2$). Then $f$ is a valid partition \href{https://artofproblemsolving.com/wiki/index.php?title=Iff}{iff} $(Af)(v) = d(v)$, where $d$ is the degree of $v$, for all $v$ (this is taken over $F_2$). So we want all solution to $Af = d$. Note that if $Af = d, Ag = d$, then $A(f - g) = 0$, so $f - g$ is in the \href{https://artofproblemsolving.com/wiki/index.php?title=Nullspace&action=edit&redlink=1}{nullspace}. Thus in particular the number of solutions, if non-zero, is the size of the nullspace of $A$, which is $2^{dim(Null(A))}$ by considering all linear combinations of any basis of $Null(A)$ over $F_2$. Also let $i$ be such that $i(v) = 1$ for all $v$. Then clearly $Ai = 0$, so $dim(Null(A)) > 0$, establishing that the number of ways to do this is $2^k$, $k > 0$. Thus we need only prove the existence of a solution.

 Since we can add a new vertex connected to exactly one previously existing vertex without changing the problem, \href{https://artofproblemsolving.com/wiki/index.php?title=Without\_loss\_of\_generality}{without loss of generality} all vertices have odd degree. Then we want to show that $i \in Col(A)$. But it is a well-known fact in linear algebra that $Col(A) = Null(A^T)^{\perp} = Null(A)^{\perp}$ since $A$ is symmetric. Thus we need to show that if $f \in Null(A)$, $f$ is perpendicular to $i$ and we will be done. So let $f \in Null(A)$. Take the submatrix of $A$ consisting of the rows and columns $v$ such that $f(v) = 1$. Then, since $f \in Null(A)$, the sum of each row in this submatrix must be $0$ in $F_2$. Thus the total number of $1$s in this submatrix is even. But since it is symmetric, the total number of $1$s off of the diagonal is even, so the total number of $1$s on the diagonal is even. But since every vertex has odd degree, the entire diagonal of $A$ consists of $1$s, so this says that the size of the diagonal of this submatrix is even. But this is also the number of $v$ such that $f(v) = 1$, so $f(v) = 1$ for an even number of $v$, thus is perpendicular to $i$, and we have our result.

 



\textbf{Solution 2 (group theory)}

Define an \textit{order} to be a set of instructions, one instruction given to each mathematician.  Each mathematician is told to either \textit{move} or to \textit{stay} (we can think of this as stay is $0$, and move is $1$).  Now take some good configuration.  Consider the set of all orders which, when performed on this configuration, give us another good configuration.  Note the \href{https://artofproblemsolving.com/wiki/index.php?title=Identity}{identity} order, $(0, 0, \ldots, 0)$, is in this set. We claim this set is an \href{https://artofproblemsolving.com/wiki/index.php?title=Abelian\_group}{abelian group} under \href{https://artofproblemsolving.com/wiki/index.php?title=Composition}{composition}.

 \textit{Proof}: Clearly each is its own \href{https://artofproblemsolving.com/wiki/index.php?title=Inverse}{inverse}, there is the identity, and the operation is clearly associative and commutative (because it's equivalent to addition of n-dimensional vectors $\mod{2}$).  So it suffices to show this set of orders is closed under composition.

 Consider any mathematician.  If, in one of these orders, he is told to stay, then the number of his friends who are told to move must be even.  Similarly, if he is told to move, then the number of his friends who are told to stay must be even.  Now just consider two orders $A$ and $B$ and you can show that in $A \times B$ the same property will hold using \href{https://artofproblemsolving.com/wiki/index.php?title=Parity}{parity}.

 Now that we've shown it is a group (which we will call $G$), we'll prove it has order two.  Let $I$ be the identity.

 Let $\{I, T_1\} = H_1$, where $T$ is some element of $G$.  Now pick an element $T_2$ of $G$ which is not in $H_1$.  Notice that because the elements of $H_1$ are distinct, the elements of $\{T_2X|X \in H_1\}$ are distinct (if two elements of that set were the same, multiply by each on the right by $T_2^{ - 1}$ and you have a contradiction).  Now notice that for any $A, B \in H_1$, if we were to have $T_2A = B$, then $T_2 = BA^{ - 1} \in H_1$.  Therefore, $H_1$ and $\{T_2X|X \in H_1\}$ are disjoint and of the same size.  Moreover, the product of any element in the first group and any element in the second group is a member of the second group.  Therefore, these two groups together form a group of order $2^2$.  Call this $H_2$. You progressively build larger and larger subgroups of $G$ until you get to $G$ itself, whose order must then be a power of two.  Therefore, the number of good configurations of the mathematicians was a power of two.

 This, of course, was all assuming $I$ existed and was in the group.

 



\textbf{Solution 3}

Let $n$ be the number of participants at the conference. We proceed by induction on $n$.

 If $n = 1$, then we have one participant who can eat in either room; that gives us a total of $2 = 2^1$ options.

 Let $n\geq 2$. The case in which some participant, $P$, has no friends is trivial. In this case, $P$ can eat in either of the two rooms, so the total number of ways to split $n$ participants is twice as many as the number of ways to split $(n - 1)$ participants besides the participant $P$. By induction, the latter number is a power of two, $2^k$, hence the number of ways to split $n$ participants is $2\times 2^k = 2^{k+1}$, also a power of two. So we assume from here on that every participant has at least one friend.

 We consider two different cases separately: the case when some participant has an odd number of friends, and the case when each participant has an even number of friends.

 \textbf{Claim:} Some participant, $Z$, has an odd number of friends.

 Remove $Z$ from consideration and for each pair $(X,Y)$ of $Z$'s friends, reverse the relationship between $X$ and $Y$ (from friends to strangers or vice versa).

 \textbf{Claim.} The number of possible seatings is unchanged after removing $Z$ and reversing the relationship between $X$ and $Y$ in each pair $(X,Y)$ of $Z$'s friends.

 \textit{Proof of the claim.} Suppose we have an arrangement prior to $Z$'s departure. By assumption, $Z$ has an even number of friends in the room with him.

 If this number is 0, the room composition is clearly still valid after $Z$ leaves the room.

 If this number is positive, let $X$ be one of $Z$'s friends in the room with him. By assumption, person $X$ also has an even number of friends in the same room. Remove $Z$ from the room; then $X$ will have an odd number of friends left in the room, and there will be an odd number of $Z$'s friends in this room besides $X$. Reversing the relationship between $X$ and each of $Z$'s friends in this room will therefore restore the parity to even.

 The same reasoning applies to any of $Z$'s friends in the other dining room. Indeed, there will be an odd number of them in that room, hence each of them will reverse relationships with an even number of individuals in that room, preserving the parity of the number of friends present.

 Moreover, a legitimate seating without $Z$ arises from exactly one arrangement including $Z$, because in the case under consideration, only one room contains an even number of $Z$'s friends.

 \textbf{End Claim}

 Thus, we have to double the number of seatings for $(n - 1)$ participants which is, by the induction hypothesis, a power of 2. Consequently, for $n$ participants we will get again a power of 2 for the number of different arrangements.

 \textbf{Case 2:} Each participant has an even number of friends.

 In this case, each valid split of participants in two rooms gives us an even number of friends in either room.

 Let $(A,B)$ be any pair of friends. Remove this pair from consideration and for each pair $(C,D)$, where $C$ is a friend of $A$ and $D$ is a friend of $B$, change the relationship between $C$ and $D$ to the opposite; do the same if $C$ is a friend of $B$ and $D$ is a friend of $A$. Note that if $C$ and $D$ are friends of both $A$ and $B$, their relationship will be reversed twice, leaving it unchanged.

 Consider now an arbitrary participant $X$ different from $A$ and $B$ and choose one of the two dining rooms. [Note that in the case under consideration, the total number of participants is at least 3, so such a triplet $(A, B; X)$ can be chosen.] Let $A$ have $m$ friends in this room and let $B$ have $n$ friends in this room; both $m$ and $n$ are even. When the pair $(A, B)$ is removed, $X$'s relationship will be reversed with either $n$, or $m$, or $m + n - 2k$ (for $k$ the number of mutual friends of $A$ and $B$ in the chosen room), or 0 people within the chosen room (depending on whether he/she is a friend of only $A$, only $B$, both, or neither). Since $m$ and $n$ are both even, the parity of the number of $X$'s friends in that room will be therefore unchanged in any case.

 Again, a legitimate seating without $A$ and $B$ will arise from exactly one arrangement that includes the pair $(A, B)$: just add each of $A$ and $B$ to the room with an odd number of the other's friends, and then reverse all the relationships between a friend of $A$ and a friend of $B$. In this way we create a one-to-one correspondence between all possible seatings before and after the $(A,B)$ removal.

 Since the number of arrangements for $n$ participants is twice as many as that for $(n - 2)$ participants, and that number for $(n - 2)$ participants is, by the induction hypothesis, a power of 2, we get in turn a power of 2 for the number of arrangements for $n$ participants. The problem is completely solved.

 



\textbf{Solution 4 (proving only existence of solution)}

By induction on the number of mathematicians, assuming solution exists for $n\leq k$ mathematicians. For $n=k+1$  mathematicians, if each of them has even degree then we're done. Otherwise, let $v$ be a node with odd degree, and denote its neighbors by $N(v)$. We form a new graph $G'$ by deleting $v$ and flipping all edges between members of $N(v)$. That is, for each pair of $a$ and $b$ in $N(v)$, they're connected in $G'$ if and only if they're not in $G$. By induction we know a solution $S'$ exists for $G'$, and it is easily verifiable that the similar solution $S$ exists in $G$ by adding $v$ to the side in $S'$ with even number of neighbors in $N(v)$.

 <i>Alternate solutions are always welcome. If you have a different, elegant solution to this problem, please \href{https://artofproblemsolving.com/wiki/index.php?title=2008\_USAMO\_Problems/Problem\_6&action=edit}{add it} to this page.</i>

 


\end{enumerate}